On identical Swendsen--Wang data for the 2D Ising model, a small CNN and an MLP trained only far from the transition give single-size 0.5-crossing estimates of $T_c$ within a fraction of the temperature grid step at $L=32$ (CNN \rhval{main/cnn/ising/tc_error_l32/mean:3}, MLP \rhval{main/mlp/ising/tc_error_l32/mean:3}), better than a PCA susceptibility read-out (\rhval{main/pca/ising/tc_error_l32/mean:3}) and not distinguishable (Wilcoxon $p=\rhval{analysis/paired-tests/ising/conf_vs_cnn_l32_p/mean:2}$ for the CNN) from the central peak of learning by confusion, whose mean error is nominally the lowest at $L=32$ (\rhval{main/confusion-mlp/ising/tc_error_l32/mean:3}). Logistic regression on raw spins is blind to the order because of the spin-flip symmetry, and gauge fixing, the control that isolates this, recovers only part of the performance. An FSS collapse of the classifier output did not improve the estimate here, the CNN collapse exponent was biased, and the confusion collapse failed in all variants. The MLP estimate depends more on the training-set size than on the excluded window, and this is not only a matter of the number of gradient updates; the CNN estimate does not. All of this holds for a small study (three sizes, ten seeds, untuned reimplemented baselines) and should be read as a measurement protocol and a set of negative results rather than a ranking of methods.
